\lyseisE{Επανάληψη κεφαλαίου Β2}
\solr{Β2}{1} Το $\mathrm{Α}'$ είναι από την άλλη μεριά του $\xi$, πάλι $3{,}4$ cm από τον άξονα. Τα $\mathrm{Α}$ και $\mathrm{Α}'$ απέχουν $6{,}8$ cm.
\solr{Β2}{2} Ισόπλευρο: $3$. Ορθογώνιο: $2$. Ρόμβος: $2$. Κύκλος: άπειροι.
\solr{Β2}{3} $5$ cm (ισαπέχει από τα άκρα).
\solr{Β2}{4} $5$ cm.
\solr{Β2}{5} Η ευθεία που περνά από το $\mathrm{Α}$ και το μέσο της $\mathrm{ΒΓ}$ (το ύψος από το $\mathrm{Α}$).
\solr{Β2}{6} Ως προς το $0$: $2$. Ως προς το $1$: $4$ (αφού $\frac{-2+4}{2}=1$).
\solr{Β2}{7} $14$ cm.
\solr{Β2}{8} Άξονα συμμετρίας έχουν τα Α, Β, Η, Τ, Φ. Κέντρο συμμετρίας έχουν τα Η, Ν, Σ, Φ.
\solr{Β2}{9} Τέσσερις γωνίες $80^\circ$ και τέσσερις $100^\circ$.
\solr{Β2}{10} $2x+(x+60^\circ)=180^\circ$, $3x=120^\circ$, $x=40^\circ$. Οι γωνίες είναι $80^\circ$ και $100^\circ$.
\solr{Β2}{11} Όχι, γιατί οι εντός εναλλάξ δεν είναι ίσες.
\solr{Β2}{12} Της διαγωνίου $\mathrm{ΑΓ}$.
\solr{Β2}{13} Στο μέσο της διαγωνίου, $6{,}5$ από κάθε κορυφή.
\solr{Β2}{14} $\mathrm{Α}'\mathrm{Β}'=8{,}5$ cm και είναι παράλληλο προς το $\mathrm{ΑΒ}$.
\solr{Β2}{15} α) Σωστό. β) Σωστό. γ) Λάθος. δ) Σωστό.
\solr{Β2}{16} α) Σωστό. β) Λάθος. γ) Σωστό. δ) Σωστό.
